% ===================================================================
%  Section IX: Discussion
% ===================================================================

\subsection{What the System Verifies}

The current system checks a precise stack of relation statements: (i)~a
private transition belongs to a committed replay dataset via Merkle
membership; (ii)~fixed-point TD target, TD error, SmoothL1 loss, and public
claimed outputs are recomputed correctly; (iii)~Q-values are anchored to
committed MLP weights through model-grounded forward execution;
(iv)~gradients, SGD deltas, and post-update checkpoints are consistent with
the checked forward pass; (v)~multi-step training fragments chain all of the
above through deterministic sampling and checkpoint linking; and
(vi)~proof-manifest chunk-chain aggregation binds externally verified child
fragment manifests.

Each of these layers has both a Python semantic oracle implementation (Layer~2
in \Cref{fig:verification-stack}) and an SP1 proof-backed implementation
(Layer~3--5), with explicit backend coverage documented in
\Cref{tab:backend-coverage}.

\subsection{What the System Does Not Claim}

Several important limitations are inherent to the current scope:

\begin{itemize}[leftmargin=1.5em]
  \item \textbf{No full training proof.} The non-goals listed in
  \Cref{sec:threat-model} apply here unchanged.

  \item \textbf{No Adam proof.} The tiny update scope is limited to simple SGD.
  Modern DQN implementations typically use Adam; proving Adam's per-parameter
  adaptive learning rates and momentum terms remains future work.

  \item \textbf{No large-scale proofs.} The benchmarks use CartPole and
  LunarLander, with PointMaze backing the public dataset-scaling rows.
  Atari-scale or MuJoCo-scale proofs would require more efficient
  neural-network circuits, and we measured by how much. MinAtar's
  DQN~\cite{young2019minatar} carries $132{,}566$ parameters on Breakout: a
  $16$-filter $3\times3$ convolution over a $10\times10\times4$ state, then a
  $128$-unit fully connected layer. Widening our hidden layer to $132{,}565$
  parameters costs $126\times$ the cycles of the $836$-parameter network proved
  here at a matched step count --- $2.033$ billion against $16.13$ million at
  $k = 4$, both in execute mode over the committed
  \texttt{lunarlander-random} dataset. A $k = 156$ fragment costs $494.6$
  million cycles at the proved width under the same settings, which puts a leaf
  of the wider network near $62$ billion. The guest evaluates a dense Linear--ReLU--Linear
  network, so this prices the parameter count and not the convolutions.

  \item \textbf{Recursive aggregation needs a GPU.} The aggregate guest
  verifies child proofs cryptographically for $T \in \{16, 32, 64\}$ and for all
  three whole-run binary trees, but only under a CUDA prover: peak device memory
  is $18.4$~GB and flat in $T$, while the CPU prover did not complete within
  $61$~GB of host memory. The $T \in \{32, 64, 128\}$ rows remain proof-manifest
  chains by choice, so the table reports both aggregation modes side by side.

  \item \textbf{No honest collection for public data.} Minari/D4RL imports
  are source-integrity commitments. The system does not prove that the
  original data was collected honestly.

  \item \textbf{The prover declares the value bound.} The sampler seed is
  derived from the committed dataset root and the chunk's global step, so a
  prover cannot grind it for a favourable subset. The bound
  \texttt{q\_abs\_max\_fp} is a public input the relation checks for
  self-consistency, but the prover chooses its value: a verifier reads the
  declared range off the statement rather than being given a canonical one.

  \item \textbf{Guest binaries are reproducible, but only inside the pinned
  image.} Provenance records \texttt{guest\_elf\_sha256} for every proved row.
  Building the guests in the SP1 release image at a pinned tag makes that hash
  portable: five of the eight workspaces built to the same digest on an Ubuntu
  WSL host and on the EC2 instance, byte for byte. The qualification is that the
  identity belongs to the image rather than to the source --- a different SP1
  release compiles the same Rust to a different guest, and the tag is what the
  repository pins. Outside the image the hash reverts to identifying the build:
  before we pinned it, the same source under two directories produced two
  hashes. Cycle counts remain the stronger invariant, deterministic in the
  (ELF, input) pair.

  \item \textbf{Tamper rejection does not cover relation semantics.} Each of
  the 236 cases perturbs a committed field and checks that verification fails.
  A relation that consistently computes the wrong thing passes all of them: an
  earlier sampler drew the same eight transitions in every chunk of a chain, so
  a 1248-step run was 156 repetitions over eight rows, and the benchmark stayed
  green throughout.
\end{itemize}

\subsection{Why Scoped Claims Remain Valuable}

A narrow claim is not a weakness if it represents the right decomposition
point. Offline DQN training contains several distinct verifiability
sub-problems: membership in committed trajectories, TD arithmetic,
neural-network evaluation, gradient computation, optimizer semantics, and
trace composition. By proving each sub-problem independently and composing
them through checkpoint chaining, the system provides a modular foundation
that can be extended to cover additional sub-problems without redesigning
the entire pipeline.

From a systems/security perspective, the contribution should be read as a
\emph{reproducible ZK-backed verification pipeline for RL-specific relations},
not as an end-to-end zkML training system. The theorem-to-artifact mapping
(\Cref{sec:relations}) ensures that every formal claim has traceable
implementation, test, and benchmark evidence.

\subsection{Scalability Considerations}

The SP1 proof cost data (\Cref{tab:proof-cost}) reveals that proving is
prototype-scale. Under a CUDA prover the most expensive relation is Groth16
child-proof aggregation at $1{,}590$\,s; the training fragments that dominate a
run cost seconds ($3.3$\,s at $k\!=\!8$), and a proof covering an entire
1248-step run costs $192$\,s. Verification stays flat at $0.054$--$0.13$\,s
regardless, except for Groth16, whose in-circuit SNARK check takes $68$\,s.
Several scaling directions are available:

\begin{itemize}[leftmargin=1.5em]
  \item \textbf{GPU acceleration, measured.} Giving the remaining six hosts a
  CUDA path and re-proving their ten rows cut prove times by $51$--$62\times$
  ($90.4$\,s to $1.78$\,s for Fwd-TD MLP, $50.2$\,s to $0.86$\,s for Merkle
  membership) at cycle counts identical to the digit. The speedup is in the
  prover, not in the relation, which is why the table is read in cycles.
  \item \textbf{Longer runs under one proof.} The aggregate guest already
  verifies child proofs in-circuit, so trace length is bounded by budget rather
  than by construction: the 1248-step trees cost roughly one GPU hour each, and
  the same shape extends to 4992 steps for about four.
  \item \textbf{Precompile optimization.} SP1 precompiles for SHA-256 and
  matrix operations could reduce cycle counts for forward/backward MLP
  computations.
  \item \textbf{Minibatches.} The relation fixes \texttt{batch\_size} to 1,
  and this is the scaling direction that matters most for what the proof says
  rather than for how fast it runs: an ablation isolating each difference
  between the tuned and the provable configuration finds the minibatch and the
  target-sync interval, not the optimizer or the clipping rule, are what cost
  the provable configuration its return.
\end{itemize}

\subsection{Comparison with Existing ZK Training Systems}

\Cref{tab:comparison} provides a qualitative comparison with related ZK
training verification systems. Our system addresses RL-specific semantics
(replay membership, Bellman targets, target-network synchronization) and
provides an artifact-level tamper-rejection benchmark for the supported
relations.

\begin{table*}[!t]
\centering
\small
\caption{Qualitative comparison with related ZK training verification systems.
\checkmark: supported; $\sim$: partial; ---: not addressed.}
\label{tab:comparison}
\begin{tabular}{lccccccc}
\toprule
\textbf{Feature} &
\textbf{PoL~\cite{jia2021proofoflearning}} &
\textbf{zkDL~\cite{zkdl2023}} &
\textbf{Kaizen~\cite{kaizen2024}} &
\textbf{zkPoT~\cite{zkpotvicinity2025}} &
\textbf{vCNN~\cite{vcnn2024}} &
\textbf{Ours} \\
\midrule
ZK proof backend & --- & \checkmark & \checkmark & \checkmark & \checkmark & \checkmark \\
Training verification & $\sim$ & \checkmark & \checkmark & \checkmark & --- & $\sim$ \\
RL-specific relations & --- & --- & --- & --- & --- & \checkmark \\
Replay membership proof & --- & --- & --- & --- & --- & \checkmark \\
Bellman/TD target proof & --- & --- & --- & --- & --- & \checkmark \\
Target-network proof & --- & --- & --- & --- & --- & \checkmark \\
Dataset commitment & --- & --- & \checkmark & $\sim$ & --- & \checkmark \\
Sampling/order integrity & --- & --- & \checkmark & --- & --- & \checkmark \\
Checkpoint chaining & --- & --- & \checkmark & --- & --- & \checkmark \\
Multi-step composition & --- & \checkmark & \checkmark & --- & --- & \checkmark \\
Recursive aggregation & --- & --- & \checkmark & --- & --- & \checkmark \\
Tamper-rejection benchmark & --- & --- & --- & --- & --- & \checkmark \\
Artifact reproducibility & $\sim$ & --- & --- & --- & --- & \checkmark \\
\bottomrule
\end{tabular}
\vspace{2pt}
\begin{minipage}{0.96\textwidth}
\footnotesize
For our system, $\sim$ under training verification means proof-backed
training fragments and whole-run aggregation over a committed dataset; it does
not denote a proof of model selection, optimizer state, or initialization.
Recursive aggregation is marked supported because the aggregate guest verifies
each child proof in-circuit, under a CUDA prover. Kaizen is marked supported
for dataset commitment and for sampling integrity: it commits the training set
and derives each epoch's ordering publicly. The row we claim there is not the
existence of the mechanism but its fit to a sampler that draws with
replacement, which \Cref{subsec:kaizen-applied} works through.
\end{minipage}
\end{table*}

\subsection{What a Supervised Proof of Training Covers Here}
\label{subsec:kaizen-applied}

The claim that a supervised proof-of-training construction does not discharge
these obligations should be argued rather than asserted, so we take the closest
system --- Kaizen~\cite{kaizen2024}, a zkPoT for deep networks with recursive
composition --- and apply its training step to our relation line by line.

Its per-iteration statement covers the forward pass, the backward pass and the
parameter update for a batch drawn from a committed dataset, with the residual
at the output layer formed from the \emph{labels carried by that batch}. Three
things follow, and none of them is a deficiency of Kaizen in its own setting.

\begin{itemize}[leftmargin=1.5em]
  \item \textbf{No line produces a label.} The residual reads a label that the
  commitment already fixes. A TD target does not exist until the target network
  evaluates the next state, so an offline-DQN statement must contain a second
  forward pass whose inputs are a checkpoint, not data, and must bind that
  checkpoint to the chain. There is nothing in a supervised iteration to
  specialize for this; the computation is simply absent.
  \item \textbf{No line advances a second model.} Target-network
  synchronization is a step-indexed event that rewrites the network the labels
  come from. A construction with one model has no place to put it, and a
  verifier reading a correct-looking supervised transcript cannot tell whether
  the sync schedule it assumed was the one the prover ran.
  \item \textbf{The permutation solves order, not selection.} A public
  permutation makes every record appear exactly once per epoch, so grinding the
  seed permutes a set the verifier already knows. Sampling with replacement
  from a buffer that a run never exhausts leaves the \emph{subset} to the
  prover, and it is the subset, not the order, that decides what the reported
  return is a statement about.
\end{itemize}

\noindent Composition transfers; semantics do not. Kaizen's recursive structure
is exactly the shape our aggregation needs, and we use the same idea. What does
not transfer is the content of the statement being composed, which is why the
obligations in \Cref{sec:relations} are stated over checkpoints, sync events
and a commitment-derived sampler rather than over batches and labels.

\subsection{Where the Prover's Time Actually Goes}
\label{subsec:cost-shape}

One cost result in this paper is worth stating as a property of the problem
rather than of our implementation, because a second, independent system reaches
it from the opposite direction. Recomputing the shares in Kaizen's own prover
breakdown gives the following, for a $10.1$M-parameter VGG-11 and a
$61.7$k-parameter LeNet, where ``commitment'' is the cost of committing
to model state rather than of proving any arithmetic on it:

\begin{center}\footnotesize
\begin{tabular}{@{}lrr@{}}
\toprule
\textbf{Prover time component} & \textbf{VGG-11} & \textbf{LeNet} \\
\midrule
Commitment generation & $51.3\%$ & $55.1\%$ \\
Proof of aggregation & $16.1\%$ & $10.3\%$ \\
Proof of the verifier circuit & $12.0\%$ & $27.3\%$ \\
Proof of gradient descent itself & $20.7\%$ & $7.3\%$ \\
\bottomrule
\end{tabular}
\end{center}

\medskip
\noindent Committing to the model dominates, and it dominates \emph{more} as
the network gets smaller: on the $61.7$k-parameter network the learning step
the system exists to prove accounts for $7.3\%$ of the prover's time. Our own
relation showed the same shape from a different construction. Hashing model
state four times per training step cost $76\%$ of the cycles in a $[4,64,2]$
leaf; replacing three of those hashes with a structural equality against the
previous step's output --- equal models have equal commitments, and a few
hundred integer comparisons are cheaper than a SHA-256 over a rendered model
--- made a real $k=156$ leaf $1.79\times$ cheaper and the whole $15$-node tree
$1.45\times$, at byte-identical public outputs.

Two very different stacks, a proof system built for the task and a
general-purpose zkVM, therefore agree that at the model sizes anyone can
currently prove, the
binding to model state, not the arithmetic of learning, is the thing to
optimize. We read this as a statement about the current cost regime of
proof-of-training rather than about either system, and it is the reason we
report cycles per relation rather than seconds per epoch.
